\section{Scope and limitations}
\label{sec:limits}

\textbf{Design.} Twenty games are five seeds with four seat rotations of one roster. No seed and rotation pair was repeated, and each state was reached once. The games cannot separate a system from its seat, its opponents or its luck. Per-system counts describe what each system encountered: Opus faced 58 financing episodes and Grok 38. We make no ranking.

\textbf{Text.} The tier labels and the 61-of-74 reading are model readings without human validation. Number matches are exact but see only figures that were written. Opus's text comes from a different channel than the others' because of the refusal condition, and Grok's applied calls sometimes disagree with its content. Every statement about what an agent wrote should be read with its channel.

\textbf{Mechanism.} We ran no intervention on prompt, menu, rules, memory or response handling. The corpus does not identify why two inventories stopped at the listed-only total, whether a written assessment shaped the action or followed it, or what causes Grok's channel conflicts.

\textbf{Reference points.} The continued-financing result is an argument from source, not a measured baseline. Three payment witnesses and one menu search come from an earlier engine audit that we did not rerun. We computed no least-cost liquidation and no trade values, so nothing here is a statement about optimal play.

\textbf{Selection.} The necessary-collateral partition, the bargaining contrast and the incidents were chosen after looking at the games.

\section{Conclusion}

When a debt arrived, these agents usually wrote down what their assets could raise, and the figure was usually right to the dollar, with no price list and no rules in the prompt. Their first actions mostly followed the asset bound, and six of the nine departures were one system's calls contradicting its own text. Continuing to finance was enough to pay every payable debt, so the informative behavior lies in the first assessments and in the choice of what to give up. The same accounting appears in bargaining, where agents argue from a rival's building budget and from who would inherit an estate. Why a count stopped short, and whether a tactic paid, need a different design.
